\documentclass[11pt,a4paper]{article}

\usepackage{amsmath}
\usepackage{amssymb}
\usepackage{siunitx}
\usepackage{booktabs}
\usepackage{geometry}
\usepackage{hyperref}
\usepackage{xcolor}
\usepackage{framed}
\usepackage{enumitem}
\geometry{margin=2.5cm}

\newcommand{\Sf}{\mathbf{S}_f}
\newcommand{\uvec}{\mathbf{U}}
\newcommand{\dt}{\Delta t}
\newcommand{\rhocont}{\rho_{\mathrm{cont}}}

\title{From Conservation Laws to the \texttt{TampinesSteamArray} Algorithm\\
\large A derivation of the 1-D compressible PIMPLE solver, its HEM closure,\\
and the three fixes that make the Edwards blowdown plateau appear}
\author{Outram Park / TAMPINES}
\date{}

\begin{document}
\maketitle

\begin{framed}
\noindent\textbf{AI-generated document.} This derivation was generated with AI
assistance from the conservation laws and from a reading of the
\texttt{TampinesSteamArray} source. It is \emph{explanatory background}, not a
novel contribution: the finite-volume and PIMPLE material is standard textbook
content (see the references), and is reproduced here only so that the
project-specific parts at the end rest on a stated foundation.

\noindent It is written to be \emph{re-derivable}. Every non-obvious step states
\emph{why} it is taken, because a derivation you cannot reconstruct is a
derivation you do not have.

\noindent Statements about what the code does are tied to
\texttt{src/openfoam\_algorithms/rhoPimpleFoam/}. Where the implementation
departs from stock \texttt{rhoPimpleFoam}, this is flagged explicitly in
\S\ref{sec:deviations} rather than quietly idealised.
\end{framed}

\tableofcontents

%======================================================================
\section{The problem and the plan}

We want a 1-D compressible pipe solver for water/steam that can survive a
violent depressurisation (the Edwards--O'Brien blowdown), in which subcooled
liquid flashes to a two-phase mixture and chokes at the break.

The chain of reasoning is:

\begin{enumerate}[leftmargin=*]
  \item Write the conservation laws (\S\ref{sec:governing}).
  \item Discretise them by the finite-volume method, obtaining one algebraic
        equation per cell (\S\ref{sec:fv}).
  \item Observe that pressure has \emph{no equation of its own}, and manufacture
        one from continuity (\S\ref{sec:pressure}).
  \item Assemble those pieces into the PIMPLE loop (\S\ref{sec:pimple}).
  \item Close the system with IAPWS-IF97 steam tables instead of an ideal gas,
        which is where the difficulty lives (\S\ref{sec:closure}).
  \item Repair the three failures that the steam-table closure exposes:
        energy conservation (\S\ref{sec:rhocont}), compressibility
        (\S\ref{sec:psi}), and shock capturing with its stability taper
        (\S\ref{sec:knp}).
\end{enumerate}

%======================================================================
\section{Governing equations}\label{sec:governing}

For a compressible single-component fluid, with $\rho$ density, $\uvec$
velocity, $p$ pressure and $h$ specific static enthalpy:

\begin{align}
  \frac{\partial \rho}{\partial t} + \nabla\!\cdot\!(\rho \uvec) &= 0
  \label{eq:mass}\\[2pt]
  \frac{\partial (\rho \uvec)}{\partial t} + \nabla\!\cdot\!(\rho \uvec \uvec)
    &= -\nabla p + \nabla\!\cdot\!(\mu \nabla \uvec)
  \label{eq:mom}\\[2pt]
  \frac{\partial (\rho h)}{\partial t} + \nabla\!\cdot\!(\rho \uvec h)
    &= \frac{\partial p}{\partial t} + \nabla\!\cdot\!(\alpha_h \nabla h) + Q
  \label{eq:energy}
\end{align}

closed by an equation of state. For an ideal gas one writes $\rho = p/RT$; here
we instead use the IAPWS-IF97 backward flash
\begin{equation}
  \rho = \rho(p,h), \qquad T = T(p,h), \qquad \mu = \mu(p,h),
  \qquad \alpha_h = \lambda(p,h)/c_p(p,h).
  \label{eq:eos}
\end{equation}

\paragraph{Unknowns versus equations --- the orphan.} Count them: the unknowns
are $\uvec$, $p$, $\rho$, $h$. Momentum naturally yields $\uvec$; energy yields
$h$; the equation of state yields $\rho$. \emph{Nothing yields $p$.} Pressure
appears in \eqref{eq:mom} only as a gradient, and momentum is already committed
to producing $\uvec$; one vector equation cannot deliver two fields.

What is left over is continuity, \eqref{eq:mass}. So pressure is determined
\emph{implicitly}, by the requirement that the velocity field conserve mass.
This single observation dictates the entire algorithm, and we return to it in
\S\ref{sec:pressure}.

%======================================================================
\section{Finite-volume discretisation}\label{sec:fv}

\subsection{Integration and Gauss's theorem}

Integrate each law over a control volume $V_c$ and convert divergences to
surface integrals:
\begin{equation}
  \int_{V_c} \nabla\!\cdot\!\mathbf{F}\, \mathrm{d}V
  = \oint_{\partial V_c} \mathbf{F}\cdot \mathrm{d}\mathbf{S}
  \;\approx\; \sum_f \mathbf{F}_f \cdot \Sf ,
\end{equation}
where $\Sf$ is the \emph{outward} face-area vector. The outward orientation is
carried entirely by $\Sf$; it must not be applied a second time when
approximating gradients.

Continuity \eqref{eq:mass} becomes, with mass flux $\phi_f = (\rho\uvec)_f\cdot\Sf$
in \si{kg/s},
\begin{equation}
  V_c\,\frac{\rho_c - \rho_c^{o}}{\dt} \;+\; \sum_f \phi_f \;=\; 0 .
  \label{eq:disc-mass}
\end{equation}
Both terms are \si{kg/s}; omitting $V_c$ or the dot product with $\Sf$ leaves
the equation dimensionally inconsistent.

\subsection{The diffusion stencil}

For a diffusive flux the face gradient is a directional derivative in the
outward direction, i.e.\ ``ahead minus behind'' with the neighbour ahead:
\begin{equation}
  (\nabla \psi)_f \approx \frac{\psi_N - \psi_P}{\delta_f},
  \qquad
  \sum_f \Gamma_f |\Sf| \frac{\psi_N - \psi_P}{\delta_f}.
  \label{eq:diffusion}
\end{equation}
The sign is fixed by physics, independent of convention: if the neighbour is
hotter or faster, cell $P$ must \emph{gain}. Expanding \eqref{eq:diffusion}
gives one term in $\psi_P$ and one in $\psi_N$.

\subsection{The matrix that results}

Every discretised transport equation collapses to one algebraic relation per
cell, linking the cell to its neighbours:
\begin{equation}
  a_P \psi_P + \sum_N a_N \psi_N = b .
  \label{eq:row}
\end{equation}
Two distinctions must be kept apart, and conflating them is the classic error:

\begin{center}
\begin{tabular}{@{}ll@{}}
\toprule
\textbf{axis} & \textbf{question} \\
\midrule
implicit vs.\ explicit & is the value at the new time (unknown) or known? \\
diagonal vs.\ off-diagonal & is it \emph{this} cell's value or a \emph{neighbour}'s? \\
\bottomrule
\end{tabular}
\end{center}

\noindent $a_P$ is the coefficient on the cell's \emph{own} unknown --- the
matrix diagonal. The $a_N$ are off-diagonal. Both may be implicit. Deleting the
off-diagonals would leave a purely diagonal matrix, invertible cell-by-cell,
in which nothing diffuses: \emph{the off-diagonal entries are the physical
coupling}.

\subsection{The \texorpdfstring{$H$}{H} operator}

Split the matrix $A = D + N$ into diagonal and off-diagonal parts. Then
\begin{equation}
  (D+N)\uvec = b
  \;\Longrightarrow\;
  D\uvec = b - N\uvec
  \;\Longrightarrow\;
  \uvec = D^{-1}(b - N\uvec).
  \label{eq:splitting}
\end{equation}
Writing $a_P \equiv D$ and $H(\uvec) \equiv b - N\uvec$, and separating the
pressure gradient out of $b$:
\begin{equation}
  \boxed{\; a_P \uvec_P = H(\uvec) - \nabla p \;}
  \label{eq:aph}
\end{equation}
So $H$ is not an arbitrary bucket: it is precisely ``the right-hand side after
moving the off-diagonal action across.'' It contains the neighbours'
contributions (still unknowns, at the new time) \emph{plus} genuinely known
sources such as the old-time term. It is written $H(\uvec)$ because it depends
on the velocity field and must be re-evaluated as the solution updates.

Equation \eqref{eq:splitting} is a Jacobi/Gauss--Seidel splitting: the unknown
appears on both sides. \emph{That is why the scheme must iterate}, and it is the
same structural reason the outer corrector loop exists.

\subsection{Convection and its linearisation}

Convection $\nabla\!\cdot\!(\rho\uvec\uvec)$ is nonlinear --- velocity
transports itself. It is linearised by lagging the mass flux at the previous
iterate while keeping the transported velocity implicit,
\begin{equation}
  \sum_f \underbrace{\phi_f}_{\text{lagged coefficient}} \; \uvec_f^{\,\text{new}},
\end{equation}
which is a \emph{Picard} (successive-substitution) linearisation. With
first-order upwind, the face value is taken from the \emph{upstream} cell:
\begin{equation}
  \uvec_f =
  \begin{cases}
    \uvec_P, & \phi_f > 0 \quad\text{(outflow)} \;\Rightarrow\; \text{diagonal},\\
    \uvec_N, & \phi_f < 0 \quad\text{(inflow)} \;\Rightarrow\; \text{off-diagonal}.
  \end{cases}
\end{equation}
Note that the lagged quantity is a \emph{coefficient}, never the unknown: the
velocity itself remains implicit.

%======================================================================
\section{Manufacturing the pressure equation}\label{sec:pressure}

Divide \eqref{eq:aph} by $a_P$ and define
\begin{equation}
  r_{AU} \equiv \frac{1}{a_P},
  \qquad
  \mathrm{HbyA} \equiv \frac{H(\uvec)}{a_P},
  \qquad\text{so}\qquad
  \uvec_P = \mathrm{HbyA} - r_{AU}\,\nabla p .
  \label{eq:hbya}
\end{equation}
$\mathrm{HbyA}$ is the velocity a cell would have from its neighbours and its
own history alone --- everything \emph{except} the pressure force.

Now substitute \eqref{eq:hbya} into discrete continuity \eqref{eq:disc-mass}.
The only remaining unknown is $p$, and because $\nabla p$ sits inside a
divergence we obtain a variable-coefficient Laplacian:
\begin{equation}
  \sum_f (\rho\, r_{AU})_f (\nabla p)_f \cdot \Sf
  \;=\;
  V_c \frac{\rho - \rho^o}{\dt} + \sum_f \rho_f\, \mathrm{HbyA}_f\!\cdot\!\Sf .
  \label{eq:peqn-generic}
\end{equation}
This is a \emph{scalar} equation, matching the scalar unknown $p$. The orphan
now has a home.

\paragraph{Reading the right-hand side.} The term
$\sum_f \rho_f \mathrm{HbyA}_f\!\cdot\!\Sf$ is the mass imbalance of the
pressure-free velocity. The pressure equation is therefore literally: \emph{the
size of the continuity violation tells you how to correct the pressure.}

\paragraph{Compressibility enters here.} $\rho$ on the right still depends on
$p$. Linearising with
\begin{equation}
  \psi \equiv \left.\frac{\partial \rho}{\partial p}\right|_{h},
  \qquad
  \rho \approx \rho^{*} + \psi\,(p - p^{*}),
\end{equation}
and moving the resulting term to the left adds a diagonal contribution
$\psi V_c/\dt$, making the equation Helmholtz-like rather than a pure Laplacian:
\begin{equation}
  \boxed{\;
  \psi_c \frac{V_c (p_c - p_c^{o})}{\dt}
  + \sum_f s_{c,f}\Big[\phi_{\mathrm{HbyA},f}
      - (\rho\, r_{AU})_f \,\mathrm{snGrad}(p)_f\,|\Sf|\Big] = 0 \;}
  \label{eq:peqn}
\end{equation}
In the incompressible limit $\psi \to 0$ and the classical pressure Poisson
equation is recovered. Thus $\psi$ sits \emph{in the diagonal of the pressure
equation}: it is the knob controlling how strongly pressure resists change.
This will matter in \S\ref{sec:psi}.

%======================================================================
\section{The PIMPLE loop}\label{sec:pimple}

Why a loop at all? Because the equations are coupled but we solve them
\emph{segregated}, one variable at a time. Each solve therefore uses stale
values from the others, and the only repair is to go round again. A monolithic
solver --- all unknowns in one matrix --- would need no predictor--corrector at
all, but pays with one enormous, ill-conditioned system.

Two distinct lags are being repaired simultaneously:

\begin{center}
\begin{tabular}{@{}lll@{}}
\toprule
\textbf{lag} & \textbf{what is stale} & \textbf{example} \\
\midrule
nonlinearity (within an equation) & a coefficient depending on the unknown & lagged $\phi_f$ \\
coupling (between equations) & a value owned by another equation & $\nabla p$ in momentum \\
\bottomrule
\end{tabular}
\end{center}

\noindent The canonical ordering is forced, not chosen: momentum must be
assembled first (it alone supplies the $\uvec$--$p$ relation), then continuity
manufactures $p$, then velocity is corrected.

\begin{enumerate}[leftmargin=*]
  \item \textbf{Momentum predictor.} Solve \eqref{eq:aph} with the guessed
        $p^{*}$; the resulting $\uvec^{*}$ does \emph{not} conserve mass.
  \item Form $r_{AU}$ and $\mathrm{HbyA}$.
  \item \textbf{Pressure corrector.} Solve \eqref{eq:peqn}.
  \item \textbf{Velocity/flux correction} with the new $p$; continuity now holds.
  \item Repeat 3--4 (PISO inner correctors); refresh and repeat all (PIMPLE outer).
\end{enumerate}

SIMPLE, PISO and PIMPLE differ only in how many times one goes round and what is
refreshed. PIMPLE contains both as limits: one outer corrector reduces it to
PISO; dropping the time derivative and relaxing heavily gives SIMPLE.

%======================================================================
\section{Closure by steam tables, and why it breaks things}\label{sec:closure}

Replacing the ideal gas by the IAPWS-IF97 flash \eqref{eq:eos} is not a cosmetic
substitution. Inside the saturation dome:

\begin{itemize}[leftmargin=*]
  \item $\rho$ responds violently to $p$ (a small pressure drop flashes liquid to
        vapour), so $\psi$ is orders of magnitude larger than any single-phase value;
  \item the equilibrium sound speed collapses, so the flow becomes near-sonic at
        modest velocities;
  \item the flash has hard validity limits --- a state pushed below
        \SI{273.15}{\kelvin} or above \SI{1073.15}{\kelvin} does not merely lose
        accuracy, it \emph{panics}.
\end{itemize}

Each of the three fixes below is a direct consequence of one of these facts.

%======================================================================
\section{Fix 1 --- conservative energy and the continuity density}\label{sec:rhocont}

\subsection{Why the naive form leaks}

In continuous calculus the conservative and non-conservative energy forms agree:
\begin{align}
  \frac{\partial(\rho h)}{\partial t} + \nabla\!\cdot\!(\rho\uvec h)
  &= \rho\frac{\partial h}{\partial t} + h\frac{\partial \rho}{\partial t}
   + \rho\uvec\!\cdot\!\nabla h + h\nabla\!\cdot\!(\rho\uvec) \notag\\
  &= \rho \frac{Dh}{Dt}
   + h\underbrace{\Big[\frac{\partial \rho}{\partial t} + \nabla\!\cdot\!(\rho\uvec)\Big]}_{\text{continuity}} .
  \label{eq:expand}
\end{align}
The bracket vanishes \emph{only because continuity holds exactly}. If continuity
is satisfied to within a residual $r$, the two forms differ by exactly $h\,r$ ---
a spurious energy source or sink.

Discretely $r$ is generally non-zero, because the thermodynamic density
$\rho(p,h)$ returned by the flash and the density implied by the discrete mass
balance are computed by \emph{different parts of the solver} and need not agree.
During flashing, $\rho$ swings enormously and the leak $h\,r$ becomes large:
enthalpy is over-drained, the liquid is driven subcooled, and the pressure falls
straight past $p_{\mathrm{sat}}$ instead of holding the flashing plateau.

\subsection{The continuity-consistent density}

Rearranging discrete continuity \eqref{eq:disc-mass} for the new density defines
\begin{equation}
  \boxed{\;\rhocont \;\equiv\; \rho^{o} - \frac{\dt}{V_c}\sum_f \phi_f
  \;=\; \rho^{o} - \dt\,\nabla\!\cdot\!\phi \;}
  \label{eq:rhocont}
\end{equation}
By construction $(\rhocont - \rho^{o})/\dt = -\nabla\!\cdot\!\phi$ holds
\emph{exactly}, whatever the equation of state returns. Using $\rhocont$ in the
storage term therefore makes the cancellation in \eqref{eq:expand} exact,
term for term, and the energy equation reduces to the material derivative
$\rho\,Dh/Dt = \partial p/\partial t$. That reversible $\mathrm{d}h \approx
\mathrm{d}p/\rho$ is what keeps the state on the saturation dome as $p$ falls ---
i.e.\ \emph{the plateau}.

\paragraph{As implemented.} The energy time derivative is assembled by
\texttt{fvm::ddt\_coeff\_old}, giving
\begin{equation}
  \frac{\rhocont\, h - \rho^{o} h^{o}}{\dt} V_c ,
\end{equation}
with two coupled points: the old-time coefficient is the \emph{old} density
$\rho^o$ (restoring the missing $h^o(\rho-\rho^o)/\dt$ term), and the new-time
coefficient is $\rhocont$ recomputed from the \emph{final} mass flux --- not the
EOS density that the thermodynamic update wrote during the inner correctors.
Using the EOS density instead leaves a residual that spuriously \emph{over-heats}
cells, flashing them into Region~5.

$\rhocont$ is floored at \SI{1e-4}{kg/m^3}. In a clamped cell the exact
cancellation formally no longer holds --- a caveat that returns in
\S\ref{sec:knp}.

%======================================================================
\section{Fix 2 --- the compressibility \texorpdfstring{$\psi$}{psi}}\label{sec:psi}

With energy conserved, the flash initiates at the correct pressure but
\emph{overshoots} the plateau. The fault is now in the pressure equation.

A quality-frozen isothermal compressibility $\psi = \rho\,\kappa_T$ omits the
flashing contribution $(v_g - v_f)\,\mathrm{d}x/\mathrm{d}p$ and is roughly two
orders of magnitude too small inside the dome. Since $\psi$ occupies the
diagonal of \eqref{eq:peqn}, too small a $\psi$ makes the solver behave as though
the fluid were nearly incompressible ($\psi \to 0$ is exactly the incompressible
limit), so it fails to resist the depressurisation and sails through the plateau.

The remedy is to evaluate the true derivative along the thermodynamically
relevant path --- constant enthalpy, since $h$ is the transported variable and a
fast blowdown is approximately adiabatic --- by central difference on the real
flash:
\begin{equation}
  \psi \;=\;
  \frac{\rho(p+\Delta p,\,h) - \rho(p-\Delta p,\,h)}{2\,\Delta p},
  \qquad
  \Delta p = \max\!\big(10^{-3} p,\ \SI{50}{\pascal}\big).
  \label{eq:psi}
\end{equation}
The bounds are clamped to the EOS validity range and the \emph{actual}
$(p_{\mathrm{hi}} - p_{\mathrm{lo}})$ is used as the denominator, so clamping
degrades the estimate to a one-sided difference rather than corrupting it. If
both bounds collapse together the isothermal value is used as a fallback, and
$\psi$ is floored at $10^{-12}$.

With \eqref{eq:psi} the plateau emerges from the thermodynamics alone --- no
clamps, no saturation clips.

\paragraph{Cost.} Two extra $(p,h)$ flashes per cell per thermodynamic update,
and the update runs once per inner corrector: $2\,n_{\text{inner}}
n_{\text{outer}} n_{\text{cells}}$ extra flashes per timestep.

%======================================================================
\section{Fix 3 --- all-Mach shock capturing and its taper}\label{sec:knp}

\subsection{Why dissipation is needed}

In this solver the energy convection term is assembled \emph{explicitly with
central interpolation}. Central differencing of convection is unconditionally
non-monotone: at the near-sonic flashing front it rings. A pressure-based
solver has almost no numerical dissipation of its own, so the oscillation is not
damped.

\subsection{KNP central-upwind flux}

Borrow the Kurganov--Noelle--Petrova central-upwind flux from density-based
schemes. With one-sided local wave-speed estimates built from the \emph{HEM
equilibrium} sound speed $c$,
\begin{equation}
  a_L = \min\big(u_{n,L}-c_L,\; u_{n,R}-c_R,\; 0\big),
  \qquad
  a_R = \max\big(u_{n,L}+c_L,\; u_{n,R}+c_R,\; 0\big),
\end{equation}
the flux of a conserved variable $w$ is
\begin{equation}
  F = \frac{a_R F_L - a_L F_R + a_L a_R (w_R - w_L)}{a_R - a_L}.
  \label{eq:knp}
\end{equation}
The final term is the dissipation. Subtracting the corresponding central flux
(the same expression with the jump term removed) isolates it, and it is applied
as a \emph{deferred correction} added to $\phi$ rather than by rebuilding the
matrix. Left and right states are MUSCL-reconstructed with a van~Leer limiter.

\subsection{Three non-obvious choices}

\paragraph{(a) Dissipate $\rho h_e$, not $\rho E$.} The conserved energy variable
is the \emph{static} enthalpy density. Dissipating total energy
$\rho E = \rho(h_e + \tfrac12|\uvec|^2) - p$ would re-inject the $-\Delta p$
pressure work; across the strong rarefaction the large $\Delta p$ then over-cools
the near-break cell through the \SI{273.15}{\kelvin} validity edge.

\paragraph{(b) Gate on $\min(\mathrm{Ma})$, not $\max$.} The blend weight is
\begin{equation}
  \beta = \mathrm{clamp}\!\left(\frac{\mathrm{Ma}_f - \mathrm{Ma}_{\mathrm{lo}}}
  {\mathrm{Ma}_{\mathrm{hi}} - \mathrm{Ma}_{\mathrm{lo}}},\,0,\,1\right),
  \qquad
  \mathrm{Ma}_f = \min(\mathrm{Ma}_O,\,\mathrm{Ma}_N),
\end{equation}
with defaults $\mathrm{Ma}_{\mathrm{lo}}=0.3$, $\mathrm{Ma}_{\mathrm{hi}}=1.0$.
At a liquid/two-phase interface the liquid's $\sim\!\SI{1400}{m/s}$ sound speed
dominates the wave speeds, so the numerical viscosity $a_L a_R/(a_R-a_L)$ is
enormous --- yet that liquid acoustic wave is genuinely \emph{low Mach} and must
not be dissipated. Taking the minimum sees the subsonic liquid side and returns
$\beta = 0$; only where \emph{both} sides are near-sonic does KNP activate.

\paragraph{(c) No separate energy source.} The continuity dissipation is folded
into $\phi$ \emph{before} the energy equation is assembled, so
$\nabla\!\cdot\!(\phi h)$ already carries it. Adding a standalone
$\Delta F_{\mathrm{ener}}$ would double-count.

\subsection{The rarefied-tail taper}

Even so, the full \SI{600}{\milli\second} run failed at $t \approx
\SI{0.18}{\second}$. As the pipe empties, $\rhocont$ falls to its floor and the
energy-equation diagonal $\rhocont V_c/\dt$ collapses; the cell becomes
ill-conditioned, and a fixed explicit correction then produces an unbounded
enthalpy update that trips the $(p,h)$ validity edge.

Since the minimum dissipated-face density during the physics of interest is
$\approx \SI{106.5}{kg/m^3}$, a density threshold cleanly separates the runaway
from the region that matters. The blend weight is tapered:
\begin{equation}
  \boxed{\;
  g(\rho_f) = \mathrm{clamp}\!\left(\frac{\rho_f - 50}{100 - 50},\,0,\,1\right),
  \qquad \rho_f = \min(\rho_O,\rho_N),
  \qquad \beta \leftarrow \beta\, g(\rho_f) \;}
\end{equation}
Below \SI{50}{kg/m^3} the dissipation is switched off entirely (pure PIMPLE,
stable); above \SI{100}{kg/m^3} it is at full weight, so the taper is
\emph{inert} in the validated region and the plateau is unchanged. The physical
justification is that in the emptying tail there is no flashing shock left to
capture.

The effective face weight is therefore
\begin{equation}
  \beta_f = \mathrm{clamp}\!\Big(\tfrac{\min(\mathrm{Ma}_O,\mathrm{Ma}_N)-0.3}{0.7},0,1\Big)
  \cdot
  \mathrm{clamp}\!\Big(\tfrac{\min(\rho_O,\rho_N)-50}{50},0,1\Big),
\end{equation}
further gated by a validity guard requiring both cells to lie within
\SIrange{300}{1050}{\kelvin}.

%======================================================================
\section{The assembled algorithm}\label{sec:algorithm}

Per timestep, with $u^o,p^o,h^o,\rho^o$ the previous state:

\begin{enumerate}[leftmargin=*]
  \item \textbf{Outer loop}, $n_{\text{outer}}$ times:
  \begin{enumerate}
    \item \textbf{Continuity (explicit).}
          $\rho = \max(\rho^{o} - \dt\,\nabla\!\cdot\!\phi,\ 10^{-4})$.
    \item \textbf{Assemble momentum}: implicit Euler time term, first-order
          implicit upwind convection, orthogonal Laplacian diffusion; add the
          lagged hybrid momentum source if in \texttt{HybridAllMach} mode.
    \item Cache $A = \mathrm{diag}$ and $r_{AU} = V_c/A$.
    \item \textbf{Momentum predictor}: add $-V_c\nabla p$ explicitly, solve for
          $\uvec$, then \emph{remove} it again so that $H(\uvec)$ is
          pressure-free for the projection.
    \item \textbf{Inner loop}, $n_{\text{inner}}$ times:
          form $\mathrm{HbyA}$ and $\phi_{\mathrm{HbyA}}$; assemble and solve
          \eqref{eq:peqn}; under-relax and bound $p$; correct the internal-face
          fluxes $\phi = \phi_{\mathrm{HbyA}} - (\rho r_{AU})_f
          \mathrm{snGrad}(p)|\Sf|$; correct
          $\uvec = \mathrm{HbyA} - r_{AU}\nabla p$; re-close the thermodynamics
          via the $(p,h)$ flash (which also refreshes $\psi$ by \eqref{eq:psi}).
    \item \textbf{Hybrid dissipation} (if enabled): fold $\Delta\phi$ into
          $\phi$ immediately; store the momentum source for the \emph{next}
          outer corrector.
    \item \textbf{Energy}: build $\rhocont$ from the final $\phi$, assemble
          the conservative $\mathrm{ddt}$ with $(\rhocont,\rho^{o})$, explicit
          central convection, implicit $\alpha_h$ diffusion, $\partial p/\partial t$
          source, plus lateral conductance and volumetric heat sources; solve for $h$.
  \end{enumerate}
  \item Clear the per-step lateral/heat registrations.
\end{enumerate}

\paragraph{Boundary conditions.} A closed end is a Dirichlet $\uvec = 0$. The
break is a \emph{choked-flow outlet}: the critical pressure and mass flux are
obtained from the HEM critical-flow dispatcher using the last cell's $(p_0,h_0)$,
converted to an equivalent full-face velocity so that $\dot m$ is conserved
through the smaller break area, and imposed as a Dirichlet velocity (lagged one
step). Both pressure patches remain zero-gradient: a blowdown has no Dirichlet
pressure anywhere --- mass depletion sets $p$.

%======================================================================
\section{Deviations from stock \texttt{rhoPimpleFoam}}\label{sec:deviations}

These are deliberate or known, and are recorded so the document is not read as
describing textbook behaviour.

\begin{enumerate}[leftmargin=*]
  \item \textbf{Equation order.} Stock is
        \texttt{rhoEqn} $\to$ \texttt{UEqn} $\to$ \texttt{EEqn} $\to$ \texttt{pEqn};
        here energy is solved \emph{after} the pressure loop, so $h$ lags $p$ by
        one outer corrector.
  \item \texttt{rhoEqn} is recomputed every outer corrector, not only on the
        first iteration.
  \item \textbf{Energy convection is explicit and central}, not implicit limited
        upwind. This is the dispersion the KNP hybrid exists to control.
  \item No \texttt{ddtCorr} (Rhie--Chow transient flux correction) in
        $\phi_{\mathrm{HbyA}}$.
  \item No kinetic-energy or gravity-work terms in the energy equation, and no
        viscous dissipation.
  \item The \emph{momentum} time derivative is non-conservative
        ($\rho^{\text{new}}(\uvec-\uvec^{o})/\dt$) while the \emph{energy} time
        derivative is conservative. Asymmetric by design.
  \item No residual-based PIMPLE convergence test: corrector counts are fixed and
        solver performance is discarded. Only the linear-solver tolerances apply.
  \item Hybrid \emph{momentum} dissipation is lagged one outer corrector, so with
        $n_{\text{outer}} = 1$ it is computed and then discarded. The Edwards
        runs use $n_{\text{outer}} = 4$.
  \item Hybrid dissipation and the flux correction touch internal faces only.
  \item Numerical floors ($\rho,\rhocont \ge 10^{-4}$, $\psi \ge 10^{-12}$,
        $c \ge \SI{1}{m/s}$) formally break the exact discrete-continuity
        cancellation in clamped cells.
  \item Enthalpy boundary conditions are not re-stamped after a solve, so a
        Dirichlet enthalpy BC survives one step and then silently degrades to
        zero-gradient. This does not affect Edwards, which sets no enthalpy BC.
  \item No turbulence model, wall function, wall friction, mesh motion, LTS, or
        transonic \texttt{pcEqn} variant.
\end{enumerate}

%======================================================================
\section{Summary}

The three project-specific fixes are one idea in three costumes: \emph{keep the
discrete system self-consistent}.

\begin{center}
\begin{tabular}{@{}lll@{}}
\toprule
\textbf{fix} & \textbf{failure it removes} & \textbf{mechanism} \\
\midrule
$\rhocont$ + conservative ddt & plateau collapses to \SI{17.4}{psia} & exact discrete-continuity cancellation \\
$\psi = \partial\rho/\partial p|_h$ & plateau overshoots & correct stiffness in the pEqn diagonal \\
KNP + $\min(\mathrm{Ma})$ + taper & ringing; tail crash at \SI{0.18}{\second} & selective, gracefully deactivated dissipation \\
\bottomrule
\end{tabular}
\end{center}

\noindent None of them is a tuned fudge: each is the consequence of requiring
that the discrete equations mean what the continuous ones mean.

%======================================================================
\section*{References}

\begin{itemize}[leftmargin=*]
  \item H.~Jasak, \emph{Error Analysis and Estimation for the Finite Volume
        Method with Applications to Fluid Flows}, PhD thesis, Imperial College,
        London, 1996. (Finite-volume discretisation, deferred correction, the
        $a_P$/$H$ formulation.)
  \item H.~Jasak, A.~Jemcov, \v{Z}.~Tukovi\'{c}, ``OpenFOAM: A C++ Library for
        Complex Physics Simulations,'' \emph{Int.\ Workshop on Coupled Methods in
        Numerical Dynamics}, Dubrovnik, 2007.
  \item J.~H.~Ferziger, M.~Peri\'{c}, \emph{Computational Methods for Fluid
        Dynamics}, Springer.
  \item A.~Kurganov, S.~Noelle, G.~Petrova, ``Semidiscrete central-upwind schemes
        for hyperbolic conservation laws and Hamilton--Jacobi equations,''
        \emph{SIAM J.\ Sci.\ Comput.} \textbf{23}(3), 2001.
  \item W.~Wagner, H.-J.~Kretzschmar, \emph{International Steam Tables}, Springer.
  \item P.~Saha, \emph{A Review of Two-Phase Steam-Water Critical Flow Models
        with Emphasis on Thermal Nonequilibrium}, NUREG/CR-0417, BNL-NUREG-50907,
        1978.
\end{itemize}

\end{document}
